\chapter[Evolución del esquema ACM2YAML]{Evolución del esquema\\ ACM2YAML}
\label{cap:evolucion}

Este capítulo presenta la contribución C-02 (Sección~\ref{sec:introduccion:contribuciones}) y responde a la RQ2: si ACM2YAML es suficientemente expresivo para representar la capacidad contractual de las APIs de correo, y qué extensiones requiere.
Al escribir las ocho Datasheets del catálogo contra ACM2YAML v0.3 aparecieron treinta y cinco hechos del dominio que el esquema no puede expresar o que PRIME4API no consume.
El primero salió del piloto: Postmark vende seis volúmenes contratados dentro de un mismo plan, de 10.000 a 1.500.000 mensajes, y una cuota de v0.3 solo admite una cifra y un sobrecoste.
El esquema es de Gavira Sánchez y Ruiz López \citep{gavira2026tfm,ruizlopez2026tfm}, y este trabajo lo extiende con propuestas sobre sus campos, cada una con el caso real que la motiva.

El capítulo se organiza del siguiente modo.
En primer lugar, la Sección~\ref{sec:esquema:metodo} describe cómo un elemento no soportado llega a propuesta de extensión y distingue dos clases de carencia, la de expresividad y el defecto de contrato.
A continuación, la Sección~\ref{sec:esquema:inventario} presenta el inventario de los treinta y cinco elementos, con su recurrencia en las ocho APIs y el origen de las nueve propuestas.
Después, la Sección~\ref{sec:esquema:expresividad} formula las siete propuestas de expresividad, y la Sección~\ref{sec:esquema:contrato}, las dos de contrato y una rectificación de la guía.
La Sección~\ref{sec:esquema:fuera} recoge los elementos promovidos sin propuesta y la compatibilidad de las propuestas con v0.3.
Por último, la Sección~\ref{sec:esquema:respuesta} responde a la RQ2.

\section{Cómo una carencia llega a propuesta}
\label{sec:esquema:metodo}\label{sec:propuesta:conformidad}\label{sec:evolucion:metodo}\label{sec:evolucion:clases}

Las ocho Datasheets del catálogo se escriben contra la guía de ACM2YAML v0.3, el documento que describe el esquema con ejemplos, sin inventar campos \citep{apisconsumptionmodel}.
Las dos formas que quedan fuera de la guía, la unidad \texttt{MBs} de los topes de tamaño y el registro \texttt{not\_found} de lo que el proveedor no publica, ya se describieron junto al consumo del catálogo por PRIME4API (Sección~\ref{sec:propuesta:prime4api}).
Escribir sin campos propios tiene una ventaja: el catálogo sirve aunque no se acepte ninguna extensión, porque PRIME4API lo lee tal cual, con los fallos que recoge la Tabla~\ref{tab:verificacion}.

Cuando un proveedor publica algo que el esquema no puede representar, no se añade a la Datasheet ningún campo propio.
El hecho se registra como «elemento no soportado»: un hecho del dominio que ACM2YAML v0.3 no puede expresar, como los seis volúmenes contratados de Postmark, o que expresa de una forma que PRIME4API, el motor de referencia, no consume.
Cada elemento lleva un número F-nn, el caso real que lo motiva, la consecuencia para el análisis y, cuando lo hay, un esbozo de resolución.
En el repositorio de datos estos registros se llaman «fricciones» (fichero \texttt{fricciones.md}), y de ahí viene la F del número.

Un elemento no soportado se promueve, es decir, pasa a candidato de extensión del esquema, por una de dos vías.
La primera es la recurrencia: el elemento reaparece en al menos dos APIs del catálogo.
Que Postmark venda seis volúmenes contratados para su cuota mensual podría ser una rareza suya, y verlo en un segundo proveedor indica que pertenece al dominio.

La segunda vía es la traza, que admite una sola API si se cumplen a la vez tres condiciones.
La condición L1 pide que el hecho esté citado con su literal o medido en un sondeo, con la evidencia archivada en el repositorio.
La condición L2 pide la consecuencia ejecutada sobre el motor: con la Datasheet sin manipular, PRIME4API devuelve un valor erróneo o deja sin responder lo que respondería si el esquema pudiera expresar el hecho.
Pide además un caso de control: la misma operación del motor responde bien cuando el dato cabe en el esquema.
La condición L3 pide que ningún mecanismo existente de la guía cubra la carencia y que no la explique un defecto ya conocido del motor, lo que descarta una mala lectura de la guía.

Quedar promovido no equivale a ser propuesta de extensión.
Solo lleva número de propuesta, P-nn, el elemento promovido cuya forma de extensión se ha comprobado contra los campos que define la guía.
La propuesta P-05 es la excepción: no sale de ningún elemento no soportado y trata el plan cuyo precio existe pero no se publica.

Un elemento que aparece en una sola API y no deja traza sobre el motor no se añade al esquema, aunque esté bien documentado.
Queda en el inventario sin promover, y su fila anota la condición que no cumple.
El criterio mide la generalidad dentro del catálogo, y no se ha comprobado si los once elementos no promovidos (Sección~\ref{sec:esquema:inventario}) aparecen en APIs de fuera del catálogo.

La recurrencia muestra que una carencia pertenece al dominio, pero no que la forma elegida para sortearla funcione.
Esa segunda comprobación exige una herramienta que lea la Datasheet y diga si la forma elegida se procesa.
Ninguno de los repositorios del ecosistema, ni el de la guía ni los de PRIME4API y SPHERE, incluye un esquema JSON, un analizador estático o una integración continua que compruebe un fichero \texttt{datasheet.yaml}.
PRIME4API es un motor de cálculo y no un validador: lee los campos del YAML sin comprobarlos antes, y sus esquemas describen las peticiones de su propia API REST \citep{prime4apirepo}.
El repositorio de la guía aporta la norma y los ejemplos, pero no un comprobador \citep{apisconsumptionmodel}.

Gavira Sánchez describe así el editor de Datasheets de SPHERE: «Cada cambio en el documento se valida contra la sintaxis del modelo» \citep{gavira2026tfm}.
En el repositorio de SPHERE, en su versión del 10 de junio de 2026, la página de subida exige la extensión \texttt{.yaml} o \texttt{.yml} y lee el YAML solo para sacar su nombre y su versión, y el editor comprueba que el YAML se pueda leer.
En ese código no se ha encontrado una validación contra los campos del modelo.
Qué comprueba el editor en la versión desplegada de la plataforma \citep{sphere} no se ha verificado.

Por eso la conformidad se establece por uso: una herramienta de este trabajo ingiere cada Datasheet en PRIME4API y la somete a diecisiete comprobaciones fijas y a una más por cada plan (Sección~\ref{sec:propuesta:prime4api}).
Cada propuesta declara así dos cosas separadas: en cuántas APIs del catálogo aparece la carencia y si hay traza de ejecución que muestre su efecto sobre el motor.
La primera es un argumento sobre el dominio, y la segunda, sobre el esquema.

Los elementos no soportados se reparten en dos clases, y cada una se prueba de forma distinta.
Una «carencia de expresividad» aparece cuando el proveedor publica algo que el esquema no tiene dónde poner.
Postmark vende seis volúmenes contratados para su cuota mensual, cada uno con su precio y su sobrecoste.
Una cuota de v0.3 solo tiene un \texttt{value} escalar y un \texttt{overage\_cost} único, así que no hay campo donde escribir la lista.
Esta carencia se detecta al modelar, porque el fichero no admite el dato.

Un «defecto de contrato» aparece cuando el esquema sí ofrece una forma, la Datasheet la usa bien y PRIME4API falla al leerla: se rompe el contrato implícito por el que el motor lee todo lo que la guía admite.
La guía admite dejar vacío un valor cuando no hay cifra que escribir.
En Postmark y en Mailchimp Transactional queda vacío el tope de destinatarios, que la fuente no pone, y en Resend el \textit{rate} del plan Enterprise, que existe y no se publica.
El motor convierte ese valor a número sin comprobar si está vacío, y la Datasheet, aun siendo conforme, deja sin responder diez comprobaciones en Postmark y en Mailchimp Transactional (Tabla~\ref{tab:verificacion}).
Este defecto solo se detecta al ejecutar el motor sobre el fichero, porque al modelar la Datasheet se escribe sin resistencia.

La primera consecuencia de la distinción está en la vía de promoción.
La recurrencia muestra que la carencia pertenece al dominio.
La traza muestra que el motor da un resultado erróneo o deja sin responder la operación, y esa prueba vale en una sola API para cualquiera de las dos clases.
Un defecto de contrato se descubre ejecutando el motor, así que llega ya con esa ejecución.
Una carencia de expresividad sin traza, como los volúmenes de Postmark, necesita aparecer en una segunda API.

La clase de carencia decide también dónde se arregla y en qué repositorio se comunica.
Los volúmenes de Postmark solo caben si cambia el esquema y la cuota admite una lista de volúmenes.
El campo vacío admite el arreglo en el esquema, en el motor o en los dos, y cada propuesta dice en cuál.
El cambio de la cuota se comunica en el repositorio de la guía, y la comprobación del campo vacío, en el del motor.
La separación se ha obtenido sobre un solo esquema, ACM2YAML v0.3, y un solo motor, PRIME4API, por lo que no se afirma que valga fuera de ese caso.

\section{Inventario}\label{sec:esquema:inventario}\label{sec:evolucion:inventario}

El piloto con Postmark dejó ocho elementos no soportados por ACM2YAML v0.3, los de la Tabla~\ref{tab:fricciones}.
Las siete APIs siguientes abrieron treinta números, de F-09 a F-38, y tres se retiraron o se absorbieron en otros (F-17, F-19 y F-20), así que el inventario cierra con treinta y cinco.
La Tabla~\ref{tab:recurrencia} recoge cada uno con el número de APIs en que aparece, si promueve a candidato de extensión y por qué vía, y la propuesta que produce cuando la hay.

El criterio de promoción de la Sección~\ref{sec:esquema:metodo} pide dos APIs o más o, en una sola, la traza con sus tres condiciones (L1 a L3).
Con ese criterio promueven veinticuatro de los treinta y cinco elementos, veintiuno por recurrencia y tres por traza.
Los otros once aparecen en una sola API y les falta L2, L3 o las dos.
Si en una API el elemento se sospecha pero no se ha comprobado, esa API no suma.
Al comprobarse solo puede sumar, así que cada recuento de la Tabla~\ref{tab:recurrencia} es un mínimo.

La mayoría de los recurrentes aparece en cuatro APIs o más: trece de veintiuno.
F-06 y F-07 aparecen en las ocho APIs: el primero es un \texttt{workload} (el campo que declara, con un mínimo y un máximo, cuánto consume una petición) acotado por una métrica distinta de la suya, y el segundo, los \textit{add-ons} que modifican la capacidad.
F-05, la dimensión \textit{Results Per Request} que v0.3 no serializa, aparece en siete.
En el otro extremo, F-12 (la recepción consume la cuota de envío) y F-14 (el plan caduca) cumplen el umbral con dos APIs justas.
Lo mismo ocurre con F-34 (dos topes de tamaño sobre la misma llamada) y F-37 (un tope de capacidad que fija el cliente).

\begin{table}[tbp]
\centering
\footnotesize
\setlength{\tabcolsep}{3pt}
\renewcommand{\arraystretch}{0.85}
\begin{tabular}{|c|>{\raggedright\arraybackslash}p{0.61\textwidth}|c|c|c|}
\hline
\textbf{Id} & \textbf{Elemento} & \textbf{APIs} & \textbf{Vía} & \textbf{Prop.} \\
\hline
F-01 & Un \textit{rate} que existe y cuyo valor no se publica & 4/8 & recurr. y traza & P-04 \\
\hline
F-02 & Volumen contratado con precio y sobrecoste propios & 5/8 & recurr. & P-01 \\
\hline
F-03 & Consumo multiplicativo entre dimensiones anidadas & 3/8 & recurr. & P-02 \\
\hline
F-04 & Restricciones declaradas que el proveedor admite no aplicar o que la respuesta no refleja & 4/8 & recurr. & discus. \\
\hline
F-05 & \textit{Results Per Request} documentada y no serializable & 7/8 & recurr. & --- \\
\hline
F-06 & \texttt{workload} acotado por una métrica distinta de la suya & 8/8 & recurr. & P-03 \\
\hline
F-07 & \textit{Add-ons} que modifican dimensiones de capacidad & 8/8 & recurr. & --- \\
\hline
F-08 & Los \textit{endpoints} disponibles cambian con el \textit{tier} & 4/8 & recurr. & --- \\
\hline
F-09 & El documento caduca entero, y sus datos caducan uno a uno & 5/8 & recurr. & --- \\
\hline
F-10 & Dos productos de precio con unidades de consumo distintas & 3/8 & recurr. & --- \\
\hline
F-11 & Límites de recursos por plan, que no son consumo & 6/8 & recurr. & --- \\
\hline
F-12 & La recepción consume la cuota de envío & 2/8 & recurr. & --- \\
\hline
F-13 & El sobrecoste es opcional y no existe en el plan gratuito & 5/8 & recurr. & --- \\
\hline
F-14 & El plan tiene fecha de caducidad & 2/8 & recurr. & --- \\
\hline
F-15 & Dos interfaces con topes propios sobre la misma cuota & 5/8 & recurr. & --- \\
\hline
F-16 & El sobrecoste tiene un tope que impone el proveedor & 1/8 & traza & P-07 \\
\hline
F-18 & La capacidad depende del estado de la suscripción o de la reputación, no del plan & 6/8 & recurr. & P-06 \\
\hline
F-21 & Un cupo de destinatarios distintos que se gasta al probar & 1/8 & traza & P-08 \\
\hline
F-22 & Una credencial que solo se gestiona desde la web & 4/8 & recurr. & --- \\
\hline
F-23 & Un comodín de prefijo con excepción & 1/8 & L3 & --- \\
\hline
F-24 & Una cuota superada que retiene hasta un número fijo de envíos y deja de entregar los siguientes & 1/8 & traza & P-09 \\
\hline
F-25 & Un \textit{endpoint}, dos modos de envío y una carga que depende de otro campo & 1/8 & L3 & --- \\
\hline
F-26 & Un precio que depende de la geografía y una sola divisa en el fichero & 3/8 & recurr. & --- \\
\hline
F-27 & Un exceso que cambia el plan en vez de pagarse por unidad, disparado por otra dimensión & 1/8 & L2 & --- \\
\hline
F-28 & Una lista de IP autorizadas decide si la llamada entra & 3/8 & recurr. & --- \\
\hline
F-29 & Un precio por unidad sin cuota incluida de la que colgarlo & 1/8 & L3 & --- \\
\hline
F-30 & Precio por tramos marginales sobre el mismo consumo & 1/8 & L2 y L3 & --- \\
\hline
F-31 & La unidad de gobierno es el par cuenta $\times$ región & 1/8 & L2 & --- \\
\hline
F-32 & Precio unitario degresivo sin planes que lo sostengan & 1/8 & L2 & --- \\
\hline
F-33 & Una franquicia que no caduca ni se renueva: la cuota sin ventana & 1/8 & L2 & --- \\
\hline
F-34 & Dos topes de tamaño sobre la misma llamada, uno del mensaje y otro de la petición & 2/8 & recurr. & --- \\
\hline
F-35 & El literal de un tope trae una cifra fraccionaria y el máximo de un \texttt{workload} es entero & 1/8 & L2 y L3 & --- \\
\hline
F-36 & Dos contadores sobre la misma ventana, y un solo caudal por nodo & 1/8 & L2 & --- \\
\hline
F-37 & Un tope de capacidad que fija el cliente, no el proveedor & 2/8 & recurr. & --- \\
\hline
F-38 & Dos cuotas concurrentes de la misma unidad & 1/8 & L2 & --- \\
\hline
\end{tabular}
\caption[Inventario de elementos no soportados sobre las ocho APIs del catálogo]{Inventario de elementos no soportados sobre las ocho APIs del catálogo. APIs: en cuántas de las ocho aparece el elemento entero, sumando la variante que cumple el enunciado con otro mecanismo. Vía, recurr.: dos APIs o más. traza: una sola API con las tres condiciones cumplidas (Sección~\ref{sec:esquema:metodo}). En los once no promovidos, la columna Vía indica la condición de la traza que no se cumple: L2, la consecuencia ejecutada sobre PRIME4API con caso de control, o L3, que ningún mecanismo de la guía cubra la carencia ni la explique un defecto conocido del motor. Prop., discus.: promovido que se discute sin formular extensión. Los números que faltan corresponden a elementos retirados o absorbidos por otros. La matriz completa, con la marca de cada API y su fuente, está en \texttt{03\_datos/\allowbreak RECURRENCIA\_2026-09-23.md} del repositorio del trabajo.}
\label{tab:recurrencia}
\end{table}

La Tabla~\ref{tab:propuestas} dice de dónde sale cada propuesta.
De las nueve, ocho salen de un elemento promovido.
Cinco entran por recurrencia.
Una de ellas, P-04, sale de F-01, el \textit{rate} sin valor publicado, y tiene también la traza del campo vacío en tres Datasheets.
Las otras cuatro salen de F-02 (el volumen contratado con precio propio), F-03 (el consumo multiplicativo), F-06 (el \texttt{workload} acotado por otra métrica) y F-18 (la capacidad que depende de la suscripción o de la reputación).

Tres entran por traza, y en las tres PRIME4API da un resultado que el proveedor contradice.
En F-16, el tope que SendGrid impone al sobrecoste, el motor promete un volumen que SendGrid rechaza.
En F-21, el cupo de destinatarios distintos de MailerSend, el motor declara alcanzable en cero segundos un objetivo que la medición no consiguió enviar.
En F-24, la cuota superada de Brevo que retiene un número fijo de envíos, el motor la trata como un corte estricto: al llegar a la cifra no se envía nada más.
Brevo, en cambio, publica esa cuota con una cola de reintento.

La novena, P-05 (el plan con precio que existe y no es público), no tiene fila en el inventario.
Su carencia apareció al evaluar la Datasheet de Resend con PRIME4API, que deja de evaluar el documento entero cuando el precio de un plan lleva el valor \texttt{custom}.
Se anotó, como la forma de \texttt{workload} de la Sección~\ref{sec:esquema:contrato}, entre las discrepancias de la guía con PRIME4API, y no entre los elementos no soportados, así que no tiene número F-nn.
Los otros dieciséis promovidos no llevan número de propuesta.

\begin{table}[tbp]
\centering
\small
\setlength{\tabcolsep}{3pt}
\begin{tabular}{|l|>{\raggedright\arraybackslash}p{0.30\textwidth}|c|c|l|l|}
\hline
\textbf{Propuesta} & \textbf{Qué pide} & \textbf{Elemento} & \textbf{APIs} & \textbf{Vía} & \textbf{Clase} \\
\hline
P-01 & Volúmenes contratados en la cuota, con precio y sobrecoste propios & F-02 & 5/8 & recurr. & expresividad \\
\hline
P-02 & Consumo de una petición como producto de dimensiones anidadas & F-03 & 3/8 & recurr. & expresividad \\
\hline
P-03 & Tope de una carga expresado en otra métrica & F-06 & 8/8 & recurr. & expresividad \\
\hline
P-04 & Marca de valor no publicado junto al campo & F-01 & 4/8 & recurr. y traza & contrato \\
\hline
P-05 & Plan con precio que existe y no es público & --- & --- & motor & contrato \\
\hline
P-06 & Estado de la suscripción y reputación, declarados una vez y con su disparador & F-18 & 6/8 & recurr. & expresividad \\
\hline
P-07 & Tope del sobrecoste & F-16 & 1/8 & traza & expresividad \\
\hline
P-08 & Cuota de entidades distintas, sin periodo y con ámbito de cuenta & F-21 & 1/8 & traza & expresividad \\
\hline
P-09 & Comportamiento de la cuota al superarse & F-24 & 1/8 & traza & expresividad \\
\hline
Rectif. & Forma de lista de \texttt{workload} en la guía & --- & --- & --- & guía \\
\hline
\end{tabular}
\caption[Propuestas del capítulo y su origen]{Propuestas del capítulo y su origen. Elemento: el elemento no soportado del inventario del que sale cada propuesta. APIs: su recuento sobre las ocho del catálogo. Vía, recurr.: la carencia aparece en dos APIs o más. traza: una sola API con las tres condiciones de la Sección~\ref{sec:esquema:metodo}. motor: PRIME4API falla al leer una forma que la guía admite. Clase: expresividad, si el dominio dice algo que v0.3 no puede escribir, o contrato, si v0.3 lo admite y el motor no lo consume. En P-04 los dos apoyos no se refieren al mismo hecho: la recurrencia es la de F-01, un \textit{rate} sin cifra, y la traza es la del campo vacío en Postmark, Resend y Mailchimp Transactional, que solo en Resend es un \textit{rate}. P-05 no sale de ningún elemento del inventario, y el criterio de promoción no se le aplica. La última fila rectifica la guía y no extiende el esquema.}
\label{tab:propuestas}
\end{table}

Al pasar de las cuatro primeras APIs (Postmark, Resend, SendGrid y MailerSend) a las ocho, ningún elemento que promovía deja de hacerlo, y dos que no promovían pasan a hacerlo.
F-18, la capacidad que depende de la suscripción o de la reputación, aparece ahora en seis APIs, porque en Postmark y en Resend se comprobó al revisar las ocho.
F-22, la credencial que solo se gestiona desde la web, aparece en cuatro.
Entre ellas cuenta Amazon SES, donde esa credencial se gestiona en IAM, el servicio de identidades de AWS, fuera de la API de correo, una variante que cumple el enunciado con otro mecanismo y por eso suma.

Dos reglas fijan cómo se cuenta.
La primera es que una aparición parcial no suma.
Una variante suma si llega por otro mecanismo al mismo efecto sobre quien integra la API, como la credencial de Amazon SES gestionada en IAM.
No suma si cumple solo una parte del enunciado, como el volumen contratado de Brevo, que no lleva sobrecoste.

F-01 aparece en las ocho APIs, pero solo cuenta en las cuatro que admiten que el límite existe y no le dan cifra.
Postmark describe el error \texttt{429} sin umbral \citep{postmarkapi}, y Amazon SES escribe que su caudal de envío «varies based on your specific use case» \citep{sesdocs}.
Mailchimp no cifra la cuota horaria de la cuenta de pago\footnote{Registro de procedencia de la cuota horaria en \texttt{03\_datos/\allowbreak apis/\allowbreak mailchimp/\allowbreak datasheet.yaml} del repositorio del trabajo, que remite a \url{https://mailchimp.com/developer/transactional/docs/reputation-rejections/}.}.
Mailgun solo publica, en sus preguntas frecuentes, un caudal de 1.500 peticiones por segundo por IP de cliente, y deja sin cifra el de la cuenta\footnote{\url{https://www.mailgun.com/ai-info/}. La cita literal está en el registro de procedencia de \texttt{03\_datos/apis/mailgun/datasheet.yaml}.}.

En las otras cuatro, Resend, SendGrid, MailerSend y Brevo, una cifra cubre una parte del límite y deja otra sin cubrir.
Brevo, por ejemplo, cifra los \textit{endpoints} de su API REST y no su relé SMTP.
Contar esas cuatro como plenas habría subido el recuento de F-01 de cuatro a ocho con situaciones que no son la misma.

La segunda regla afecta a F-08, los \textit{endpoints} que cambian con el \textit{tier}.
El registro de cada API lo enuncia como el plan real que combina \textit{tier} y volumen contratado, pero el volumen contratado ya lo cubre F-02, y en F-08 solo cuenta la superficie de la API.
Solo cuenta que el plan cambie los \textit{endpoints} disponibles, como el envío por lotes, que la Datasheet de MailerSend omite de su plan gratuito.
Los límites de recursos por plan se cuentan aparte, como F-11.
Con esa regla, F-08 aparece en Postmark, SendGrid, MailerSend y Brevo.
F-08 no produce propuesta, porque v0.3 ya admite esa forma: un \textit{endpoint} que un plan no ofrece se omite de ese plan.
Promueve igualmente, porque la recurrencia cuenta en cuántas APIs aparece el patrón y no comprueba si el esquema ya lo expresa.

En treinta y una combinaciones de elemento y API la aparición no está verificada.
En cuatro filas bastaría con confirmar una sola de esas combinaciones para cambiar la vía.
F-31 (los límites que rigen por cuenta y región), F-35 (el tope fraccionario frente al máximo entero) y F-36 (dos contadores sobre la misma ventana) pasarían a recurrentes, y F-16 pasaría de la traza a la recurrencia.
En las demás filas con apariciones sin verificar la vía no cambia, porque ya promueven por recurrencia.

F-38, dos cuotas concurrentes de la misma unidad, cuenta en Resend y pasaría también a recurrente si se confirma que Mailgun aplica las dos a la vez.
Esa combinación no cuenta entre las treinta y una: está en suspenso, porque Mailgun publica cada cuota en una página distinta y ninguna de las dos dice si se aplican a la vez.

F-37 es el caso más frágil.
Aparece entero en Mailgun y promueve gracias a Brevo, cuya solución Corporate deja a la cuenta maestra fijar cuántos correos puede enviar cada una de sus subcuentas, una variante que cumple el enunciado con otro mecanismo.
Si esas variantes no se contaran, F-37 se quedaría solo en Mailgun.

\section{Propuestas de expresividad}
\label{sec:esquema:expresividad}\label{sec:evolucion:expresividad}

Siete de las nueve propuestas responden a una carencia de expresividad: el proveedor publica un límite y ACM2YAML v0.3 no tiene dónde escribirlo.
Son la cuota con varios volúmenes contratados (P-01), el consumo por petición que multiplica dos factores (P-02) o que tiene la cota en otra métrica (P-03) y los estados de la suscripción y de la reputación (P-06).
Completan las siete el tope del sobrecoste (P-07), el cupo de destinatarios distintos que no se repone (P-08) y lo que hace la cuota con lo que la supera (P-09).
Las dos restantes, P-04 y P-05, son defectos de contrato, porque v0.3 admite la forma y el motor no la consume.
Son un límite que existe y cuyo valor no se publica y un plan con precio que existe y no es público.
Las extensiones se escriben contra la guía de v0.3 \citep{apisconsumptionmodel} y no sustituyen ningún campo que ya exista.

\subsection{P-01. La cuota declara volúmenes contratados}
\label{subsec:evolucion:p01}

\textbf{El caso.} Los tres planes de pago de Postmark se anuncian como \textit{«Starting at 10,000 emails/month»} \citep{postmarkpricing}.
La estructura completa solo se ve desde el panel de una cuenta propia: seis volúmenes contratados, de 10.000 mensajes por 15 USD a 1.500.000 por 775 USD, con negociación directa por encima del último.
En el \textit{tier} Basic el sobrecoste baja de 1,80 a 1,70 USD por millar al pasar de 125.000 a 300.000 mensajes (Figura~\ref{fig:tiervolumen}, que cruza los tres \textit{tiers} con los seis volúmenes).

Resend vende la misma estructura y la publica en su página de precios \citep{resendpricing}.
Su plan Pro tiene dos volúmenes, de 50.000 mensajes por 20 USD y de 100.000 por 35 USD.
Su plan Scale recorre seis, de 100.000 mensajes por 90 USD a 2.500.000 por 1.150 USD, y el sobrecoste baja de 0,90 a 0,46 USD por millar.
Resend permite demostrar la carencia con una página pública de precios, sin cuenta propia.
Mailgun publica su lista de volúmenes en un selector de la página de precios, y en su plan Scale el sobrecoste baja de 1,10 a 0,40 USD por millar \citep{mailgunpricing}.

La carencia aparece en cinco de las ocho APIs: Postmark, Resend, Mailgun, SendGrid y MailerSend (Tabla~\ref{tab:recurrencia}).
Brevo también vende volúmenes contratados, pero sin sobrecoste, y esa variante cumple solo una parte del enunciado, así que no suma al recuento (Sección~\ref{sec:esquema:inventario}).

\textbf{Qué permite v0.3.} Una cuota de v0.3 declara un \texttt{value} escalar y un solo sobrecoste en \texttt{overage\_cost}.
No cabe la lista de volúmenes contratados ni el sobrecoste propio de cada uno.
Las Datasheets de Postmark, Resend y Mailgun modelan por eso el volumen mínimo de cada plan.
La lista completa solo queda en el registro de procedencia, en prosa que ninguna herramienta lee.

\textbf{Consecuencia.} La Datasheet de Resend afirma que Scale son 100.000 mensajes con un sobrecoste de 0,90 USD por millar.
Es cierto en el primer volumen contratado y falso en el último, donde el mismo plan son 2.500.000 mensajes a 0,46 USD por millar.
Comparar ese plan con el de un proveedor de volumen fijo usando el único número serializado falsea la comparación de precios.

\textbf{Propuesta.} Que una cuota admita una lista de volúmenes contratados, cada uno con su precio y su sobrecoste, más una entrada explícita para el tramo negociado.
El Extracto de código~\ref{lst:p01} muestra la forma sobre el plan Scale de Resend.

\needspace{18\baselineskip}
\begin{lstlisting}[style=datasheetyaml,caption={P-01: volúmenes contratados dentro de un plan, con precio y sobrecoste propios (caso Resend, plan Scale, tres de los seis volúmenes).},label={lst:p01}]
quota_email_transactional_monthly_scale:
  unit: emails
  period: {value: 1, unit: MONTH}
  window_type: fixed_billing_cycle
  contracted_volumes:
    - value: 100000
      price: 90
      overage_cost: {price: 0.90, value: 1000}
    - value: 500000
      price: 350
      overage_cost: {price: 0.70, value: 1000}
    - value: 2500000
      price: 1150
      overage_cost: {price: 0.46, value: 1000}
\end{lstlisting}

\textbf{Efecto sobre PRIME4API.} El motor calcula la curva de coste con el precio del plan y el sobrecoste de la cuota.
Con un único par de valores la curva es una recta, y el resultado solo es correcto en el volumen modelado.
Con la lista, la curva cambia de pendiente en cada volumen contratado, que es lo que el proveedor cobra.
La recomendación de plan que ya ejecuta el TFM de Gavira Sánchez (su Tabla 5.12, para 100.000 correos con sobrecoste) depende de esa curva \citep{gavira2026tfm}.

\subsection{P-02 y P-03. El \texttt{workload} deja de ser una terna escalar}
\label{subsec:evolucion:p02}

Las dos propuestas salen de dos elementos no soportados que recurren por separado: F-03, el consumo de una petición multiplica dos factores, y F-06, la cota de la petición está en otra métrica.
Los dos piden que el consumo de una petición deje de expresarse con una sola terna \texttt{unit}, \texttt{min}, \texttt{max}.
Se presentan juntas porque reformar el campo por separado obligaría a proponer dos veces la misma forma de \texttt{workload}.

\textbf{El caso de P-02.} El \textit{endpoint} \texttt{/emails/batch} de Resend acepta \textit{«up to 100 batch emails at once»} \citep{resenddocs}.
Cada mensaje admite cincuenta destinatarios, y un sondeo con cuenta propia contra Resend mostró que el tope cuenta la suma de \texttt{to} y \texttt{cc} y no cada campo por separado.
Contar cada destinatario como un correo podría ser una elección de modelado, pero Resend declara cómo cuenta los destinatarios: \textit{«Multiple To, CC, or BCC recipients count as separate emails»} \citep{resenddocs}.
El consumo de una llamada es, por tanto, la suma de los destinatarios de sus mensajes, y su máximo, el producto de los dos topes: 100 mensajes por 50 destinatarios.

En Postmark, \texttt{/email/batch} admite quinientos mensajes por llamada y cincuenta destinatarios en total por mensaje \citep{postmarkapi}.
La documentación declaraba ese tope por campo, y el sondeo del piloto (Sección~\ref{sec:propuesta:piloto}), que midió el tope sobre la suma de \texttt{to} y \texttt{cc}, lo corrigió.
Postmark no declara cómo cuenta ese consumo, y en su caso el producto de los dos factores se deduce de los dos topes medidos.

Amazon SES es el tercer caso: su envío en lote recibe una lista de entradas con sus propios destinatarios, y su cuota cuenta destinatarios y no mensajes \citep{sesdocs}.
SendGrid, MailerSend y Brevo ponen un tope sobre el total de destinatarios de la petición, y esa variante acota el total sin multiplicar dos factores, así que cumple solo una parte del enunciado y no suma al recuento (Sección~\ref{sec:esquema:inventario}).

\textbf{Consecuencia de P-02.} La Datasheet de Resend serializa \texttt{max: 100}, que es el número de mensajes y no el de unidades consumidas.
El factor de consumo por petición (el CRF, que la Datasheet serializa en \texttt{workload}) alimenta las preguntas del motor sobre agotamiento de cuota y sobrecoste.
Con ese valor, el motor subestima hasta cincuenta veces el consumo de una petición, y quien consulta el motor no puede distinguir ese número, que es plausible, de uno correcto.

\textbf{El caso de P-03.} El \textit{endpoint} \texttt{/email/bulk} de Postmark no tiene tope de mensajes.
El proveedor admite \textit{«as many recipients as you need, as long as your total payload (including attachments) stays under 50 MB»} \citep{postmarkapi}.
La cota efectiva de la llamada es un tamaño en megabytes y no un número de correos.
Resend también mide en bytes, y después de codificar, el tamaño de cada mensaje: \textit{«max 40MB per email, after Base64 encoding of the attachments»} \citep{resenddocs}.
En los dos casos el contenido útil que cabe depende de la codificación.
F-06 aparece en las ocho APIs del catálogo, con el tamaño expresado en MB, en MiB o en MB por segundo según el proveedor.

\textbf{Consecuencia de P-03.} En Postmark el campo \texttt{max} de \texttt{/email/bulk} queda sin cifra, porque el tope real está en megabytes, y en Resend el tope de 40 MB solo consta en el registro de procedencia.
Las otras seis Datasheets declaran el tope de tamaño como un segundo \texttt{workload} del \textit{endpoint} con unidad \texttt{MBs}, una forma que la guía no recoge (Sección~\ref{sec:propuesta:prime4api}).

\textbf{Propuesta.} Que \texttt{workload} admita dimensiones anidadas que se multiplican, y que \texttt{max} pueda expresarse en una métrica distinta de la unidad declarada.
El Extracto de código~\ref{lst:p02} muestra la primera forma sobre Resend y la segunda sobre Postmark.

\begin{lstlisting}[style=datasheetyaml,caption={P-02 y P-03: composición de dimensiones anidadas y cota expresada en otra métrica (caso Resend para P-02 y caso Postmark para P-03).},label={lst:p02}]
# P-02. El consumo es el producto de los dos factores.
workload:
  - unit: messages
    min: 1
    max: 100
    composes:
      unit: emails
      max_per_unit: 50
# P-03. El tope real esta en otra metrica: el tamano de la peticion.
workload:
  - unit: emails
    min: 1
    max_by: {unit: MBs, value: 50, scope: payload}
\end{lstlisting}

\textbf{Efecto sobre PRIME4API.} Las dos formas alimentan el mismo cálculo.
El factor de consumo por petición deja de ser un entero y pasa a depender de la petición, de modo que el motor tiene que evaluar el consumo máximo y no solo el declarado.
El cambio es mayor que el de P-01, porque afecta también a las operaciones que recorren la Datasheet, como la que devuelve los rangos del factor de consumo de cada \textit{endpoint}, además de a las que calculan capacidad y coste.

\subsection{P-06. El estado de la suscripción y la reputación se declaran una vez}
\label{subsec:evolucion:p06}

\textbf{El caso.} MailerSend publica cuatro estados de reputación con nombre (\textit{Healthy}, \textit{Under review}, \textit{Rebuilding} y \textit{Sending paused}) y da cifra para dos de ellos \citep{mailersenddocs}.
Cuando la cuenta entra bajo revisión, \texttt{/v1/email} pasa de 120 a 10 peticiones por minuto y el envío en lote, \texttt{/v1/bulk-email}, de 10 a 1 petición por minuto, sin cambiar de plan.
Un sondeo sobre una cuenta recién creada, que nunca se aprobó ni contrató plan, midió esos mismos 10 y 1 peticiones por minuto, una sola dirección en \texttt{to} por envío y cinco mensajes por lote.
Esa cuenta no es la que describe el valor por defecto de la Datasheet, que corresponde a la cuenta aprobada.

El caso obliga a separar tres nociones.
La cuenta es la identidad del cliente ante el proveedor, y puede tener varias suscripciones: Resend muestra bajo una misma cuenta una suscripción transaccional y otra de \textit{marketing}.
El «estado de la suscripción» es la fase en que el proveedor sitúa esa suscripción, como la aprobación pendiente de MailerSend, el entorno de pruebas de Amazon SES o el modo de prueba de Mailchimp Transactional.
La «reputación» es la valoración que hace el proveedor de cómo consume el cliente, por ejemplo por los errores que provocan sus peticiones, y amplía o reduce la capacidad sin cambiar de suscripción.

En MailerSend, la cuenta sin aprobar es un estado de la suscripción, y la revisión que baja \texttt{/v1/email} a 10 peticiones por minuto es un efecto de la reputación.
La carencia, contada sobre los estados que cambian los topes de una suscripción, recurre en seis de las ocho APIs (Tabla~\ref{tab:recurrencia}): además de MailerSend, en Postmark, Resend, Amazon SES, Mailchimp Transactional y Mailgun.

\textbf{Qué permite v0.3.} La guía documenta alias dentro de un \textit{endpoint} (bloques con nombre que llevan los límites de cada modo de consumo, aquí de cada estado) y usa MailerSend como ejemplo \citep{apisconsumptionmodel}.
Su vocabulario de ejemplo tiene dos nombres, \texttt{healthy\_reputation} y \texttt{under\_review\_reputation}.
La Datasheet de MailerSend adopta ese alias de estado que v0.3 ya ofrece y añade un tercero, \texttt{unapproved\_account}, para la cuenta sin aprobar, que el proveedor documenta y el ejemplo de la guía no recoge.
El estado sano conserva el valor por defecto, y el motor devuelve una fila por estado donde antes devolvía una sola.
Con ese alias, v0.3 ya escribe los topes de cada estado de MailerSend, y lo que falta es lo que el alias no cubre.

\textbf{Qué no cubre.} Al aplicar el mecanismo a MailerSend aparecen tres carencias.

La primera es que el alias no tiene campo para el disparador del cambio de estado.
Un \texttt{trigger} escalar dentro del alias se ignora sin aviso.
Escrito al nivel del \textit{endpoint}, el motor lo lee como un estado más, porque trata como estado todo diccionario cuya clave no sea una de las cuatro que conoce (\texttt{rate}, \texttt{quota}, \texttt{workload} y \texttt{cost\_per\_request}).
Una copia de la Datasheet con el campo dentro del alias y otra con el campo al nivel del \textit{endpoint} lo muestran.
En el comodín \texttt{/*}, que fija el caudal general de las demás rutas y lleva su \texttt{rate} al nivel del \textit{endpoint}, el motor devuelve 60 peticiones por minuto sin estado y, con el campo, las devuelve bajo un estado llamado \texttt{trigger}, que no existe.

Ninguna de las tres comprobaciones de cierre de una Datasheet detecta el estado \texttt{trigger}.
El validador de formato de este trabajo, que comprueba los tipos de cada campo y las referencias entre bloques, da 16 comprobaciones superadas y 2 avisos.
El verificador de citas comprueba que cada cita literal del registro de procedencia sigue en la copia archivada de su página, cuando la hay.
La verificación por uso con PRIME4API da 20 aciertos y 1 error.
Las tres dan el mismo resultado que sin el campo.

La segunda carencia está en la guía.
Un alias admite también un \texttt{workload} propio, y PRIME4API lo lee, pero la guía solo documenta el \textit{rate} y la cuota.
En MailerSend el tope de destinatarios por petición cambia con la aprobación de la cuenta, y ese cambio cabe en un alias aunque la guía no lo diga.
Mientras la guía no lo documente, una Datasheet que la siga al pie de la letra deja fuera del alias el tope de destinatarios de la cuenta sin aprobar de MailerSend.

La tercera es que un estado de la suscripción se declara siete veces.
El estado no depende del plan ni del \textit{endpoint}, y el esquema no tiene dónde declararlo una sola vez para toda la Datasheet.
En MailerSend la referencia se repite en los siete bloques de \textit{endpoint} y nada liga las copias, de modo que cambiar el estado obliga a tocar siete sitios.
Esta tercera carencia recurre, porque la Datasheet de Amazon SES repite su alias de entorno de pruebas en ocho bloques y la de Mailchimp Transactional repite el de la cuenta de prueba en tres.

El recuento del semáforo de confianza, que etiqueta cada valor según el respaldo de su fuente, cuenta cada copia como una evidencia.
En MailerSend, el tope de 50 MB del lote aparece tres veces en verde, una por plan, copiado del \textit{endpoint} en el estado de cuenta sin aprobar.
Solo una nota del registro de procedencia impide leer esas tres copias como tres evidencias.

\textbf{Propuesta.} Se proponen tres cambios sobre el mecanismo de alias de la guía, que se conserva.
El primero es un atributo de disparador con sitio propio, que diga si el cambio de estado lo impone el proveedor o lo elige el cliente y lleve su condición en prosa.
El segundo es documentar el \texttt{workload} por alias, que PRIME4API ya admite.
El tercero es declarar cada estado una sola vez para toda la Datasheet, con las referencias de cada \textit{endpoint} apuntando a esa declaración.
El Extracto de código~\ref{lst:p06} muestra los tres sobre MailerSend.

\needspace{15\baselineskip}
\begin{lstlisting}[style=datasheetyaml,caption={P-06: estados de la suscripción y de la reputación declarados una vez, con su disparador y con el \texttt{workload} que la guía no documenta (caso MailerSend).},label={lst:p06}]
# Declaracion unica, fuera de los endpoints.
account_states:
  unapproved_account:
    trigger: {mode: provider_managed, condition: "cuenta pendiente de aprobacion"}
    rate: rate_request_email_unapproved
    quota: quota_email_monthly_unapproved
    workload: workload_email_unapproved   # admitido por el motor, no documentado
# Cada endpoint referencia el estado en lugar de repetirlo.
/v1/email:
  states: [healthy_reputation, under_review_reputation, unapproved_account]
\end{lstlisting}

\textbf{La reputación como campo.} Los estados de ejemplo de la guía llevan la reputación en el nombre, pero ningún campo dice cómo se mueve ni cuánto cambia la capacidad.
Fuera del mecanismo de alias, se propone un campo de reputación que declare cuánto amplía o reduce la capacidad y qué lo dispara.
Con él, la Datasheet registraría si el proveedor publica cómo funciona la reputación de la cuenta de forma completa, con la capacidad de cada estado, y determinista, con el umbral que hace cambiar de estado.
Ninguna de las ocho APIs la publica así.
Seis dicen que la reputación mueve la capacidad sin dar a la vez la cifra de cada estado y el umbral, y Postmark y Brevo solo la nombran para su infraestructura de envío.\footnote{Las fuentes de cada API están en \texttt{03\_datos/\allowbreak apis/\allowbreak <api>/\allowbreak fuentes.md}, en \texttt{fricciones.md} y en \texttt{05\_salida/\allowbreak \_cap5/\allowbreak hechos\_<api>.md} del repositorio del trabajo.}

\textbf{Efecto sobre PRIME4API.} El motor ya devuelve una fila por estado, así que leer la capacidad por estado no le exige ningún cambio.
Con el disparador declarado, una consulta puede separar los estados que elige el cliente de los que impone el proveedor, que es la diferencia entre una capacidad contratable y una contingencia.
Con la declaración única, el recuento de evidencia deja de contar copias.
Hace falta además una corrección del motor y no del esquema: avisar de la clave que no reconoce y no convertirla en estado.

\subsection{P-07. El sobrecoste declara su tope}
\label{subsec:evolucion:p07}

\textbf{El caso.} Desde el 26 de marzo de 2026, SendGrid limita el envío mensual de los planes Essentials y Pro a diez veces el volumen contratado, sobrecoste incluido.
Una cuenta con el plan Essentials de 50.000 correos puede enviar 500.000 al mes, y por encima recibe el error \textit{«Maximum Credits Exceeded»}.\footnote{Artículo de soporte de SendGrid, \url{https://support.sendgrid.com/hc/en-us/articles/35466138799899-Understanding-the-Maximum-Credits-Exceeded-error}, consultado el 6 de septiembre de 2026.}
El literal cuenta \textit{requests}, que el mismo artículo define como correos procesados más descartados (\textit{«Requests = Processed + Dropped»}), así que la unidad no es la petición HTTP.
El dato aparece en un artículo de resolución de errores, y la página de precios y el artículo de sobrecostes no lo mencionan.

\textbf{Qué permite v0.3.} \texttt{overage\_cost} es un precio unitario sin límite de cantidad: la guía lo define como un objeto con precio y valor, sin ningún subcampo que acote la cantidad \citep{apisconsumptionmodel}.
La Datasheet declara los 50.000 correos incluidos y el precio de cada correo extra, y eso se lee como capacidad ilimitada a ese precio.
El tope solo cabe en el registro de procedencia, en prosa, donde ninguna herramienta lo evalúa.

\textbf{Lo que responde PRIME4API.} Pedida al motor la recomendación de presupuesto para 600.000 correos sobre ese plan, responde 734,95 USD y declara el objetivo alcanzable en 59 segundos.
El punto de control por debajo del tope, 400.000 correos, devuelve 474,95 USD con la misma forma, y el motor cruza los 500.000 sin que nada cambie.\footnote{Ejecución en \texttt{05\_salida/\_cap5/prime4api\_ejecucion\_d042/tope\_sobrecoste\_2026-10-18/} del repositorio del trabajo.}
En el caso de control, con un presupuesto de 100 USD en la consulta, el mismo motor declara el objetivo inalcanzable.
El motor declara inalcanzable un objetivo cuando el límite le llega por la consulta o por la Datasheet, así que el error de los 600.000 correos procede de que la Datasheet no puede declarar el tope.

\textbf{Propuesta.} Que \texttt{overage\_cost} admita un máximo opcional, expresado como factor sobre el volumen contratado o en unidades absolutas.
Con él, la Datasheet de SendGrid declara diez veces el plan y el motor deja de prometer un volumen que el proveedor rechaza.

\subsection{P-08. Un cupo que cuenta destinatarios distintos y no se repone}
\label{subsec:evolucion:p08}

\textbf{El caso.} La cuenta de MailerSend sin aprobar puede escribir a dos direcciones distintas, y solo a esas dos \citep{mailersenddocs}.
El cupo es de la cuenta y no de la petición. Solo lo levanta la aprobación, y la documentación no le da ventana.
El listado de destinatarios de la API (\texttt{GET /v1/recipients}) devuelve cero, así que el cupo no se puede leer por ese \textit{endpoint}.
Los sondeos lo encontraron alcanzado sin haber enviado nunca un correo, y de ahí se deduce que lo gastaron intentos que no enviaron nada.
Qué petición lo gastó no se determinó.

\textbf{Qué permite v0.3.} Una cuota tiene valor, unidad y periodo, y aquí el valor cabe, dos direcciones, pero fallan la unidad y el periodo.
La unidad son direcciones distintas, que no equivalen a correos ni a peticiones.
No hay periodo, porque el cupo no se renueva.
Y, como los sondeos encontraron el cupo gastado sin un solo correo enviado, lo consume el intento y no el envío, dos cosas que el esquema no distingue.

La guía da seis ejemplos de unidad (peticiones, correos, SMS, registros, gigabytes y \textit{tokens}), a los que sus formatos estándar suman horas y minutos \citep{apisconsumptionmodel}.
Ninguna de esas unidades cuenta entidades distintas.
Una unidad nueva, como la \texttt{MBs} del catálogo, no bastaría, porque la cuota seguiría sin decir su ámbito, su falta de periodo ni qué la consume.
La Datasheet deja el cupo en el registro de procedencia, en prosa.

\textbf{Lo que responde PRIME4API.} El motor lee la Datasheet de MailerSend y atribuye a la cuenta sin aprobar una capacidad acumulada de 110 correos en un mes, por encima incluso de los 100 correos al mes que la Datasheet le declara, cuando la capacidad medida es cero.
Un objetivo de tres correos lo declara alcanzable en cero segundos y sin coste.\footnote{Ejecución en \texttt{05\_salida/\_cap5/prime4api\_ejecucion\_d042/cupo\_destinatarios/} del repositorio del trabajo.}
La cuenta del sondeo no envió ni un correo y la cuota mensual de correos quedó intacta. Solo bajó la diaria de peticiones, de 100 a 97.
Con la misma Datasheet, el motor declara inalcanzable un objetivo de 100.000 correos por el límite de caudal, así que sabe negarse cuando el límite está escrito.
El objetivo de tres correos sería alcanzable en una cuenta recién abierta, que aún conserva sus dos direcciones, y deja de serlo en cuanto el cupo se ha gastado.

\textbf{Propuesta.} Que una cuota pueda declarar cuatro cosas que v0.3 no declara.
Son que su unidad son entidades distintas, que su ámbito es la cuenta, que no tiene periodo y qué la consume, si el intento o el resultado.
En el catálogo el cupo es solo de MailerSend: las otras siete APIs se comprobaron y no lo tienen (Tabla~\ref{tab:recurrencia}).
Por eso la propuesta se sostiene en una sola API, con la respuesta errónea que PRIME4API da sobre la cuenta sin aprobar.

\textbf{Frontera con P-06.} El cupo aparece en la cuenta sin aprobar, igual que los topes reducidos de P-06, pero pide otra cosa.
En P-06 el esquema sabe escribir el número de cada estado y no sabe decir cuál rige ni qué dispara el cambio.
En P-08 el límite no cabe ni con estado ni sin él, porque su unidad no existe en v0.3.
Se separan para que ninguna carencia cuente dos veces y cada propuesta pida una sola forma.

\subsection{P-09. La cuota declara qué hace con lo que la supera}
\label{subsec:evolucion:p09}

\textbf{El caso.} La cuota del plan gratuito de Brevo es de 300 correos al día.
Al agotarla, el correo transaccional no se corta en seco: \textit{«once you reach your daily limit, up to 1,000 additional emails are held in a retry queue. Emails beyond this queue are not delivered»} \citep{brevohelpfree}.
La cuota tiene así dos umbrales, el de la cifra publicada y el de la cola, y lo que ocurre entre los dos solo vale para el correo transaccional.
El literal no dice si el cliente recibe algún error por los correos que la cola ya no admite, y la cola no se ha observado.

\textbf{Qué permite v0.3.} De los campos de la cuota, solo \texttt{overage\_cost} habla de su borde \citep{apisconsumptionmodel}.
Sin él, la cuota se lee como un corte estricto: al llegar a la cifra no se envía nada más.
Con precio cero, se lee como un exceso que no se cobra ni se acaba.
Ninguna de las dos lecturas es la de Brevo.
La Datasheet declara los 300 correos sin sobrecoste, que en v0.3 se escribe igual que un corte estricto.

\textbf{Lo que responde PRIME4API.} Con la Datasheet sin manipular, PRIME4API devuelve para esa cuota 300 correos al día sin sobrecoste y sin rastro de la cola.\footnote{Ejecución en \texttt{03\_datos/apis/brevo/evidencias/l2\_prime4api\_2026-09-18.json} del repositorio del trabajo.}
Es la misma forma que devuelve para la prueba gratuita de SendGrid, que sí es un corte estricto, mientras que Brevo retiene hasta 1.000 correos más en la cola de reintento.
El caso de control es el plan Essentials de SendGrid: su Datasheet declara sobrecoste, y el motor lo devuelve con su precio.
El motor distingue un corte estricto de un exceso que se cobra, pero no de una cola, porque la Datasheet no tiene cómo decirlo.

\textbf{Propuesta.} Que la cuota admita un bloque \texttt{overflow} con las unidades que retiene por encima de su valor y lo que hace con ellas.
Si el bloque falta, la cuota es un corte estricto.
Si hay sobrecoste, sigue en \texttt{overage\_cost}, y el bloque nuevo solo describe la franja que no se cobra.
El Extracto de código~\ref{lst:p09} muestra la forma sobre el caso de Brevo.

\begin{lstlisting}[style=datasheetyaml,caption={P-09: una cuota que retiene lo que la supera hasta un tope fijo (caso Brevo, plan gratuito, correo transaccional).},label={lst:p09}]
quota_email_daily_free:
  value: 300
  unit: emails
  period: {value: 1, unit: DAY}
  overflow:              # P-09: lo que la cuota hace al superarse
    capacity: 1000       # envios retenidos en la cola de reintento
    behavior: queued     # en cola, y el exceso no se entrega
\end{lstlisting}

\textbf{Frontera con P-07.} El tope de SendGrid de P-07 acota lo que se cobra por encima de la cuota, y la cola de Brevo describe lo que se retiene sin cobrarlo.
Las dos propuestas tocan el mismo borde de la cuota, y reunirlas en una sola forma es una decisión que corresponde a la guía.

\section{Defectos de contrato}
\label{sec:esquema:contrato}\label{sec:evolucion:contrato}

La Tabla~\ref{tab:verificacion} (Sección~\ref{sec:propuesta:prime4api}) recoge lo que ocurre al ingerir las ocho Datasheets en PRIME4API, y sus errores tienen tres causas.
Dos son defectos de contrato, es decir, formas que la guía de ACM2YAML admite y que el motor no procesa.
Son un campo declarado y vacío en Postmark, Resend y Mailchimp Transactional, y un precio que existe y no es público en Resend, SendGrid y Brevo.

La tercera causa es el fallo del propio motor al calcular la curva acumulada de MailerSend y Amazon SES, que mezcla unidades de tiempo (Sección~\ref{sec:propuesta:prime4api}), y no produce propuesta de esquema porque ningún valor de esas Datasheets lo causa.
Mailgun termina con el mismo mensaje, pero su causa no está confirmada.
La forma de \texttt{workload} es un caso distinto, porque ninguna Datasheet del catálogo usa la variante con la que falla el motor.
La guía de ACM2YAML y una Datasheet de ejemplo de su repositorio escriben el campo de dos formas distintas, y se resuelve corrigiendo la guía.

\subsection{La forma de \texttt{workload}: mapa o lista}
\label{subsec:evolucion:dv01}

La guía del esquema documenta \texttt{workload} como un mapa \citep{apisconsumptionmodel}.
Una Datasheet de ejemplo del mismo repositorio lo serializa como una lista, y el TFM de Gavira Sánchez usa también la lista en sus extractos y en el anexo de SendGrid \citep{gavira2026tfm}.
PRIME4API tolera las dos formas en las operaciones de cálculo y solo acepta la lista en las de navegación, que recorren la estructura del documento.
Una Datasheet escrita según la guía produce entonces números correctos y falla cuando la recorren el agente conversacional HARVEY4API o la plataforma SPHERE.

La forma correcta es la lista, porque la exige el modelo de consumo.
Un mapa solo puede declarar una métrica por \textit{endpoint}, y Gavira Sánchez sostiene que un \textit{endpoint} que consume varias métricas a la vez tiene que declararlas por separado \citep{gavira2026tfm}.
Las ocho Datasheets del catálogo usan la lista.
La propuesta es una rectificación: corregir la guía para que documente la lista.
De forma opcional, el motor puede convertir el mapa en lista al leerlo, para que los ficheros escritos con la redacción actual de la guía sigan siendo navegables.

\subsection{P-04. Distinguir «sin límite» de «límite no publicado»}
\label{subsec:evolucion:p04}

\paragraph{El caso.}
Resend no publica el \textit{rate} de su plan Enterprise, y su Datasheet deja presente y vacío el campo que lo declara.
Ese campo vacío dice que el límite existe y que su valor no se conoce.
En Postmark y en Mailchimp Transactional la misma forma dice lo contrario: el proveedor no pone tope en la unidad del campo.
Si solo lee el campo vacío, PRIME4API no puede distinguir el caso de Resend del de Postmark, y P-04 propone cerrar esa ambigüedad.

En Postmark y en Mailchimp Transactional la cota que se publica está en otra métrica, los bytes de la llamada, que es la carencia que cubre P-03 (Apartado~\ref{subsec:evolucion:p02}).
Postmark no pone tope de destinatarios en \texttt{/email/bulk} y lo acota por el tamaño de la petición, 50 MB.
Mailchimp niega por escrito el tope de destinatarios en su API, «there's no recipient limitation», y limita el JSON de cada llamada a menos de 10 MB\footnote{\url{https://mailchimp.com/developer/transactional/docs/outbound-email/}, consultada el 21 de septiembre de 2026.}.
En el sondeo de Mailchimp, una llamada con 5.000 destinatarios devolvió HTTP 200 sin rechazo.

P-04 se promueve por dos vías, y cada una se refiere a un hecho distinto.
La primera es la recurrencia de F-01, un \textit{rate} que existe y no tiene cifra publicada, que cuenta en Postmark, Amazon SES, Mailchimp Transactional y Mailgun (Tabla~\ref{tab:recurrencia}).
En Mailchimp ese \textit{rate} es la cuota horaria de la cuenta de pago, que el proveedor no cifra, y no el tope de destinatarios.
La segunda es la traza: el campo vacío hace fallar al motor en tres Datasheets, aunque solo en Resend es un \textit{rate}.
En Postmark y en Mailchimp es un tope de destinatarios.
Lo que P-04 generaliza es lo que muestra la traza, un campo vacío que el motor no distingue de la ausencia de límite.

\paragraph{La traza.}
PRIME4API convierte a número el valor de esos campos sin comprobar si está vacío, y falla con un error de conversión de tipo.
En Postmark caen por esta causa diez comprobaciones: la consulta por rangos del factor de consumo por petición, tres curvas de capacidad (una es la acumulada), la curva de coste, la recomendación de plan y cuatro operaciones de cálculo.
Esas cuatro dan el tiempo mínimo para alcanzar una capacidad, la capacidad acumulada en un instante, el instante en que se agota la cuota y la espera tras agotarla.
En Resend el mismo error aparece en cinco comprobaciones, y otras ocho fallan por el precio que existe y no es público, el valor \texttt{custom} del plan Enterprise.

Como el fichero de Resend contiene los dos defectos, se generaron dos variantes, cada una con un solo defecto corregido.
Con el campo vacío relleno, la herramienta de verificación de este trabajo da trece errores, todos por el precio no público.
Con el precio no público sustituido por números, da doce, todos por el campo vacío\footnote{Las tres versiones del fichero (la original y las dos variantes) y su salida están en \texttt{03\_datos/\allowbreak apis/\allowbreak resend/\allowbreak evidencias/\allowbreak control\_motor\_causas\_2026-09-23/}.}.
Cada comprobación informa solo del primer error que encuentra.
Por eso, con los dos defectos a la vez, el campo vacío aparece en cinco y el precio no público en ocho, aunque cualquiera de los dos basta por sí solo para hacer fallar doce o trece.

En una tanda aparte de diecisiete llamadas a la API del motor con la Datasheet de Mailchimp Transactional, catorce responden HTTP 500 con el mismo error de conversión, y solo contestan las tres consultas que no calculan tiempo.
Si el campo vacío del tope se rellena en el documento ya cargado por el motor, sin tocar el fichero, el error pasa a ser una división por cero.
Esa segunda caída la causa otro campo vacío, el periodo de la franquicia de la cuenta de prueba, 500 correos gratuitos que no caducan ni se renuevan (F-33).
Con los dos campos rellenos, la misma llamada devuelve HTTP 200.
La herramienta de verificación de este trabajo, con sus propias comprobaciones, da a Mailchimp diez errores con el mismo mensaje que a Postmark.

\paragraph{Por qué el esquema no ofrece salida.}
ACM2YAML v0.3 deja dos alternativas, y las dos tienen coste.
Omitir el campo hace que la Datasheet afirme implícitamente que no hay límite.
En Resend, donde el límite existe, la omisión produciría una respuesta optimista y falsa.
En Postmark y en Mailchimp coincidiría con lo que dice el proveedor, pero PRIME4API no distinguiría esa omisión de la que haría Resend con un límite que existe.
Declarar el campo vacío, la otra alternativa, deja sin responder operaciones enteras del motor: diez comprobaciones en Postmark, cinco en Resend y catorce de diecisiete llamadas en Mailchimp.

La Datasheet de Amazon SES elige la omisión para el caudal de una cuenta con acceso de producción: no lo declara y anota la ausencia en su registro \texttt{not\_found}.
Ese registro, que usan las ocho Datasheets, dice qué falta en el documento, pero no en qué campo, y una marca junto al propio campo dejaría esa anotación donde el motor la puede leer.

\paragraph{Propuesta.}
Se propone una marca explícita de valor no publicado junto al campo.
El esquema declara la marca y el motor la respeta: omite esa dimensión y sigue evaluando el resto en lugar de abandonar el documento.
Con la marca, el límite sin cifra de Resend y la ausencia de tope de Postmark y Mailchimp dejan de escribirse igual.
Resend la lleva en su caudal de Enterprise.
Postmark y Mailchimp pueden omitir el campo, y la omisión dice lo mismo que su fuente: no hay tope de destinatarios.
La cota en bytes ya la declara la Datasheet de Mailchimp como un \texttt{workload} aparte, y la de Postmark necesita P-03.

La misma distinción alcanza al periodo de una cuota, el segundo campo vacío de Mailchimp.
Ahí, igual que en su tope de destinatarios, no falta una cifra, porque la franquicia de 500 correos de la cuenta de prueba no caduca ni se renueva (Sección~\ref{sec:propuesta:catalogo}).
La marca queda para los valores que existen y no se publican, y la cuota sin periodo es otro elemento, F-33.

\subsection{P-05. Declarar un plan cuyo precio existe y no es público}
\label{subsec:evolucion:p05}

\paragraph{El caso.}
Resend no publica el precio ni la cuota de su plan Enterprise, y la Datasheet escribe en los dos campos el valor \texttt{custom}.
Es lo que publica el proveedor, y es la forma que la guía documenta para un valor negociable \citep{apisconsumptionmodel}.

\paragraph{La traza.}
PRIME4API intenta convertir la cadena \texttt{custom} a número y falla.
El fallo no se queda en el plan Enterprise.
La evaluación recorre todos los planes del documento, así que ese único valor impide evaluar también los planes Free, Pro y Scale, que tienen precio publicado.
Con ese error fallan las tres consultas de \textit{rates}, cuotas y límites, las cuatro operaciones de cálculo y la consulta del precio del plan Enterprise, que son ocho de los trece errores de la Datasheet de Resend.
Con el campo vacío relleno, los trece pasan a esta causa (Apartado~\ref{subsec:evolucion:p04}).
SendGrid y Brevo reproducen el fallo con trece errores cada una, todos con el mismo mensaje: la cadena \texttt{custom} no se puede convertir a número.

\paragraph{Por qué importa.}
En siete de las ocho APIs del catálogo hay al menos un precio que el proveedor no publica, y la excepción es Amazon SES.
En Postmark, Resend, SendGrid, MailerSend, Brevo y Mailgun es un plan a medida, y en Mailchimp Transactional es el precio por encima de cuatro millones de correos al mes.
Resend, SendGrid y Brevo escriben ese plan con \texttt{custom}, como admite la guía, y el motor falla en las tres.
Postmark, MailerSend y Mailgun lo dejan fuera de sus planes.

\paragraph{Propuesta.}
Se proponen dos cambios separables.
En el esquema, una forma explícita de declarar un plan que tiene precio aunque no lo publique, en sustitución de la cadena \texttt{custom} que el motor no sabe leer.
En el motor, aislar el fallo por plan, de modo que un plan no evaluable devuelva su error y los demás se sigan calculando.
El segundo cambio no depende del primero, y por sí solo recupera tres de los cuatro planes de la Datasheet de Resend.

Las dos propuestas de esta sección se han medido solo contra PRIME4API.
Otro motor podría tolerar el campo vacío o la cadena \texttt{custom} sin que el esquema cambiara.

\section{Lo que queda fuera}\label{sec:esquema:fuera}\label{sec:evolucion:fuera}

De los veinticuatro elementos no soportados que superan el criterio de promoción, porque aparecen en dos APIs o más o porque una ejecución de PRIME4API muestra su efecto sobre la Datasheet, dieciséis no llevan número de propuesta.
Su patrón se repite en el catálogo, pero les falta una forma concreta contrastada con el texto de la guía de ACM2YAML.
Es un límite asumido: la guía recibe su recuento, no una forma que pueda adoptar.
Llevar número depende de tener esa forma, no de la recurrencia: F-11, los límites de recursos por plan, recurre en seis APIs y no lo lleva, y F-21, el cupo de destinatarios distintos de MailerSend, aparece en una sola API y sí lo lleva.

\phantomsection\label{subsec:evolucion:f07}F-07, los \textit{add-ons} que modifican dimensiones de capacidad, recurre en las ocho APIs.
El caso común es la IP dedicada, una dirección IP de envío reservada a un solo cliente, que se paga aparte o exige un plan mínimo.
Aun así no se propone extensión, por una razón de alcance: el esquema deja los \textit{add-ons} al canal web, y el catálogo modela solo la API REST, una de las interfaces del canal API (Sección~\ref{sec:propuesta:criterios}).
Un bloque de \textit{add-ons} en la Datasheet repetiría lo que el pricing del canal web ya modela, y enlazar la capacidad con el pricing es una línea futura que Gavira Sánchez deja en el trabajo que introduce el modelo \citep{gavira2026tfm}.
El catálogo aporta a esa línea ocho casos.

F-04 reúne los límites que el proveedor declara y que la respuesta a la petición no refleja.
Plantea si un límite que la respuesta no refleja forma parte de la capacidad contractual.
Es una pregunta sobre la definición del ACM que no se resuelve con un campo nuevo, y por eso no se convierte en propuesta.
En MailerSend, el \textit{endpoint} de lote acepta con un \texttt{202} un envío con un remitente de dominio sin verificar, que la validación diferida rechaza después.
Los dos lotes medidos terminan con cinco errores de validación y ningún correo enviado\footnote{Evidencia en \texttt{03\_datos/\allowbreak apis/\allowbreak mailersend/\allowbreak evidencias/\allowbreak probe\_sandbox\_limites\_\allowbreak 2026-09-09.json} y en la repetición del mismo sondeo, en la misma carpeta.}.
Quien mida su consumo por la respuesta cuenta cinco correos que nunca salieron.

F-04 aparece también en los sondeos de Mailchimp Transactional y de Amazon SES, que respondieron con \texttt{200} por encima del límite que declara su proveedor.
Mailchimp Transactional declara que el JSON de cada llamada debe ocupar menos de 10 MB \citep{mailchimpdocs}, y respondió con \texttt{200} a una llamada de unos 40 MB\footnote{Evidencia en \texttt{03\_datos/\allowbreak apis/\allowbreak mailchimp/\allowbreak evidencias/\allowbreak sondeo\_a\_preflight\_\allowbreak 2026-09-21.md}.}.
Amazon SES declara una petición por segundo en sus acciones de gestión \citep{sesdocs}, y respondió con \texttt{200} a doce peticiones seguidas a 5,7 por segundo (Sección~\ref{sec:propuesta:casos}).
En los dos casos la Datasheet conserva la cifra publicada, y la medida no tiene campo donde escribirse.
El registro de elementos no soportados esboza una salida: distinguir en el esquema el límite publicado del aplicado, cada uno con su fuente.

F-05, los resultados por petición que fija la documentación, es una de las dos dimensiones del ACM que ACM2YAML v0.3 deja fuera de su versión operativa.
Gavira Sánchez ya recoge completar esa serialización como línea futura \citep{gavira2026tfm}, así que proponerla repetiría esa línea.
F-08, la superficie de la API que cambia con el \textit{tier}, ya cabe en v0.3: un \textit{endpoint} que un plan no tiene se omite de ese plan (Sección~\ref{sec:esquema:inventario}).

Los doce restantes están en la Tabla~\ref{tab:sinforma}.
El registro de cada API esboza la forma de once, y la de F-12 la enuncia este trabajo, pero ninguna se ha contrastado con la guía.
F-22, la credencial que solo se gestiona desde la web, está en la frontera de lo que el esquema describe.
Su forma registra qué canal ofrece cada operación, y el esquema describe precios, cuotas y límites.

Dos de los doce tienen argumentos que no dependen de cuántas APIs los presentan.
En F-09, la caducidad de los datos, el repositorio de la guía conserva una Datasheet de MailerSend cuyos precios y sobrecostes difieren de los publicados cinco meses después \citep{apisconsumptionmodel}.
Ese mismo ejemplo define caudales de 15 y 30 peticiones por minuto que ningún plan usa, un error de modelado distinto de la caducidad.
En F-10, Resend vende un producto transaccional que se factura por correos y otro de \textit{marketing} que se limita por contactos, y la Datasheet modela el primero sin poder declarar que solo cubre ese.

\begin{table}[tbp]
\centering
\footnotesize
\setlength{\tabcolsep}{3pt}
\begin{tabular}{|l|c|>{\raggedright\arraybackslash}p{0.26\textwidth}|c|>{\raggedright\arraybackslash}p{0.36\textwidth}|}
\hline
\textbf{Dónde} & \textbf{Id} & \textbf{Carencia} & \textbf{APIs} & \textbf{Forma esbozada} \\
\hline
Documento & F-09 & El documento caduca entero, y sus datos caducan uno a uno & 5/8 & Fecha de verificación por afirmación en el registro de procedencia \\
\hline
Documento & F-10 & Dos productos de precio con unidades de consumo distintas & 3/8 & Varios catálogos por servicio, o un campo que acote qué producto se modela \\
\hline
Plan & F-11 & Límites de recursos por plan, que no son consumo & 6/8 & Bloque de recursos en el plan, con nombre, valor y unidad y sin periodo \\
\hline
Plan & F-13 & El sobrecoste es opcional y no existe en el plan gratuito & 5/8 & Modo del sobrecoste declarado en el plan \\
\hline
Plan & F-14 & El plan tiene fecha de caducidad & 2/8 & Vigencia del plan y plan al que cae al expirar \\
\hline
Precio & F-26 & Un precio que depende de la geografía y una sola divisa en el fichero & 3/8 & Precio como lista de ofertas, cada una con su divisa \\
\hline
Cuota & F-12 & La recepción consume la cuota de envío & 2/8 & Declarar que un evento entrante reenviado descuenta de la cuota de envío, con un factor que depende de la ruta \\
\hline
\texttt{workload} & F-34 & Dos topes de tamaño sobre la misma llamada, uno del mensaje y otro de la petición & 2/8 & \texttt{workload} que declara el objeto medido además de la unidad \\
\hline
Interfaz & F-15 & Dos interfaces con topes propios sobre la misma cuota & 5/8 & Interfaces declaradas, cada una con sus topes, o \texttt{workload} declarado por interfaz \\
\hline
Acceso & F-28 & Una lista de IP autorizadas decide si la llamada entra & 3/8 & Condiciones de admisión declarables en el plan o en la credencial \\
\hline
Bloque nuevo & F-37 & Un tope de capacidad que fija el cliente, no el proveedor & 2/8 & Bloque de límites que configura el cliente, fuera de la capacidad contractual \\
\hline
Operación & F-22 & Una credencial que solo se gestiona desde la web & 4/8 & Qué canal ofrece cada operación, registrado por operación \\
\hline
\end{tabular}
\caption[Promovidos sin número de propuesta y con una forma esbozada]{Promovidos sin número de propuesta y con una forma esbozada en el registro de elementos no soportados de cada API, salvo F-12, cuya forma enuncia este trabajo. Dónde: la parte del esquema que tocaría la forma. APIs: recuento sobre las ocho del catálogo. Ninguna forma se ha trabajado ni contrastado con la guía.}
\label{tab:sinforma}
\end{table}

Los once elementos sin promover aparecen en una sola API, así que solo podían entrar por la traza, y a cada uno le falta al menos una de sus tres condiciones (Tabla~\ref{tab:recurrencia}).
En F-23, el prefijo \texttt{/v3/smtp/} de la API REST de Brevo (rutas HTTP, no el relé SMTP) con dos rutas exentas, el esquema ya lo expresa si se enumeran las rutas del prefijo que no tienen fila propia.
En F-25, un \textit{endpoint} de Brevo con dos modos de envío, la guía expresa los dos modos con alias, aunque no documenta un \texttt{workload} por alias.
En F-29, un precio por unidad sin cuota de la que colgarlo, falta comprobar si la guía lo cubre.

A F-32 (precio degresivo sin planes), F-33 (franquicia sin ventana), F-36 (dos contadores sobre la misma ventana) y F-38 (dos cuotas concurrentes de la misma unidad) les falta una ejecución archivada de PRIME4API que muestre su efecto con la Datasheet sin manipular.
En F-33, la división por cero del Apartado~\ref{subsec:evolucion:p04} solo aparece tras rellenar a mano otro campo vacío.
A F-30 (precio por tramos marginales) y F-35 (tope fraccionario en un máximo entero) les faltan esa ejecución y el contraste con la guía.

F-31, los límites que rigen por cuenta y región, y F-27, el exceso de Brevo que cambia el plan en lugar de pagarse por unidad, tienen ejecución, pero ninguno de los valores que devuelve se puede dar por falso.
La de F-27 devuelve un coste plano, y con él no se distingue si Brevo deja de enviar al llegar a la cuota o pasa la cuenta al plan superior.

F-36 y F-38 piden que un mismo nodo de la Datasheet, sea un plan, un \textit{endpoint} o un alias, admita varios \textit{rates} o varias cuotas a la vez.
Es lo que pide el \textit{issue} «Soportar listas de rates y cuotas en datasheet», que Gavira Sánchez abrió en el repositorio del motor \citep{prime4apirepo}, y el motor ya lo implementa aunque la guía no lo documente.

\phantomsection\label{sec:evolucion:compatibilidad}Todas las propuestas salvo tres son retrocompatibles: añaden campos opcionales que un documento escrito contra v0.3 puede no declarar, y no cambian cómo se lee ninguno de los que ya existen.
Las dos que convierten el factor de consumo por petición en una expresión, P-02 y P-03, obligan a tocar las operaciones que lo leen.
P-05, el plan con precio no público, sustituye la cadena \texttt{custom} que usan las Datasheets de Resend, SendGrid y Brevo.
La corrección de la forma de \texttt{workload} rectifica la guía para que documente la lista que ya usan las ocho Datasheets (Sección~\ref{sec:esquema:contrato}), y queda fuera de ese recuento.
Su coste recae sobre los documentos que siguen la redacción actual y escriben un mapa.

La guía no tiene mecanismo para lo que piden las cuatro últimas propuestas: el estado de la suscripción declarado una vez y con su disparador (P-06), el tope del sobrecoste (P-07), el cupo de destinatarios distintos (P-08) y la cola tras la cuota (P-09).
Aun así, declara que el esquema es preliminar y que habrá que extender campos \citep{apisconsumptionmodel}.
En P-06, la parte que documenta el \texttt{workload} por alias no cuesta nada, porque PRIME4API ya lo lee.

Las cuatro añaden claves que v0.3 no define, y un motor que rechazara las claves desconocidas invalidaría los documentos que las declaren.
Ninguna Datasheet del catálogo las declara.
Por eso la corrección del motor que pide P-06 avisa de la clave desconocida y no la rechaza: así sigue evaluando el resto del documento y no invalida estas cuatro extensiones antes de que la guía las defina.

Antes de comunicar una carencia se comprueba si ya está reportada.
En los repositorios de la guía, de PRIME4API, de SPHERE y de las pruebas de uso de Datasheets solo hay un \textit{issue} pertinente al esquema, el de listas de \textit{rates} y cuotas, que cubre F-36 y F-38.
Fuera de ese \textit{issue}, ninguna de las carencias de P-01 a P-05 está reportada, y las de P-06 a P-09 no se han contrastado con los \textit{issues} abiertos.

\section{Respuesta a la RQ2}\label{sec:esquema:respuesta}\label{sec:evolucion:respuesta}

La RQ2 pregunta si ACM2YAML es suficientemente expresivo para representar la capacidad contractual de las APIs de correo, y qué extensiones requiere.
Con las ocho APIs del catálogo, la respuesta a la primera mitad es que no.
Las ocho Datasheets se escriben sin campos propios, pero dejan en prosa del registro de procedencia, que ninguna herramienta lee, datos como los volúmenes contratados de Postmark, Resend y Mailgun o el tope del sobrecoste de SendGrid.
En tres casos, el tope de SendGrid, el cupo de destinatarios de MailerSend y la cola de Brevo, PRIME4API devuelve un resultado que el proveedor contradice.

La respuesta a la segunda mitad son las nueve propuestas de la Tabla~\ref{tab:propuestas}: siete amplían lo que el esquema expresa y dos corrigen formas que la guía admite y PRIME4API no lee.
Ocho de ellas salen de un inventario de treinta y cinco elementos no soportados, y la novena, P-05, del fallo del motor ante un plan cuyo precio no es público.
A ellas se suma una rectificación de la guía de ACM2YAML: la guía documenta \texttt{workload} como un mapa, y una Datasheet de ejemplo de su propio repositorio lo escribe como lista (Apartado~\ref{subsec:evolucion:dv01}).

Gavira Sánchez declara que ACM2YAML v0.3 captura dos de las cuatro dimensiones de capacidad del metamodelo, el volumen (las cuotas) y la frecuencia (los \textit{rates}), y deja fuera las peticiones concurrentes y los resultados por petición \citep{gavira2026tfm}.
Los resultados por petición (\textit{Results Per Request}) están documentados en siete APIs del catálogo, y Gavira Sánchez ya recoge su serialización como línea futura \citep{gavira2026tfm}.
Las peticiones concurrentes aparecen en el catálogo en los topes del relé SMTP que reúne F-15: SendGrid admite hasta 10.000 conexiones simultáneas desde un mismo servidor, y MailerSend, cinco en sus planes Free, Hobby y Starter.
Las dos cifras quedan en el registro de procedencia.

Tres propuestas tratan lo que v0.3 sí recoge pero declara con un único valor escalar: la cifra de la cuota y el consumo de cada petición.
Postmark vende seis volúmenes contratados dentro de un mismo plan, donde el esquema admite una sola cifra de cuota (P-01, Apartado~\ref{subsec:evolucion:p01}).

Las otras dos tratan el consumo de cada petición.
Una llamada al lote de Resend admite hasta cien mensajes con hasta cincuenta destinatarios cada uno y consume hasta el producto de las dos cifras, 5.000 correos.
El esquema, en cambio, describe el consumo de una petición con una sola terna de unidad, mínimo y máximo (P-02).
El tope de \texttt{/email/bulk} de Postmark es un tamaño, 50 MB por petición, y el esquema solo admite acotar una petición en su propia unidad, que es el correo (P-03).
Los dos casos se detallan en el Apartado~\ref{subsec:evolucion:p02}.

La capacidad depende además del estado de la suscripción y de la reputación de la cuenta.
Cuando MailerSend pone una cuenta bajo revisión, su envío baja de 120 a 10 peticiones por minuto sin cambiar de plan, y ese eje aparece en seis de las ocho APIs (P-06, Apartado~\ref{subsec:evolucion:p06}).
La guía permite declarar dentro de cada \textit{endpoint} un alias por estado, con la cifra de cada uno, pero obliga a repetirlo en todos los \textit{endpoints}.
Tampoco tiene campo para el disparador que cambia el estado, como la aprobación de la cuenta en MailerSend o su revisión por reputación.

Otras dos propuestas tratan lo que ocurre al superar la cuota.
SendGrid limita el envío mensual a diez veces el volumen contratado, sobrecoste incluido, y el esquema no tiene campo para ese tope (P-07, Apartado~\ref{subsec:evolucion:p07}).
Brevo retiene hasta 1.000 correos por encima de la cuota diaria y deja de entregar los siguientes (P-09, Apartado~\ref{subsec:evolucion:p09}).
F-13, el sobrecoste opcional que no existe en el plan gratuito, se refiere también al exceso sobre la cuota en cinco APIs, y su forma queda esbozada en la Tabla~\ref{tab:sinforma}.

Al margen del exceso sobre la cuota, la cuenta de MailerSend sin aprobar solo puede enviar a dos direcciones distintas, con un cupo que no se repone (P-08, Apartado~\ref{subsec:evolucion:p08}).
P-07, P-08 y P-09 aparecen en una sola API cada una.
Se promueven porque PRIME4API devuelve sobre la Datasheet sin manipular un resultado que el proveedor contradice, y la misma operación responde bien en un caso de control en que el límite le llega escrito, en la Datasheet o en la consulta.

Las dos propuestas restantes son de contrato: el esquema admite el valor que describe al proveedor y PRIME4API falla al leerlo.
En el plan Enterprise de Resend, un límite vacío no dice si no se publica o si no existe, y el motor falla al convertirlo a número (P-04, Apartado~\ref{subsec:evolucion:p04}).
El mismo plan lleva el precio \texttt{custom}, que la guía admite y que impide al motor evaluar el documento entero (P-05, Apartado~\ref{subsec:evolucion:p05}).

Además de las nueve propuestas, la Tabla~\ref{tab:sinforma} recoge doce elementos que aparecen en dos APIs o más y cuya forma está esbozada sin contrastarla con la guía.
Otros cuatro elementos promovidos quedan fuera de las propuestas (Sección~\ref{sec:esquema:fuera}): los resultados por petición, los \textit{add-ons}, los límites que la respuesta no refleja y la superficie de la API que cambia con el plan.
Los \textit{add-ons}, como la IP dedicada, aparecen en las ocho APIs, pero el esquema los deja al canal web y el catálogo solo modela la API REST, una de las interfaces del canal API.

Tampoco hay propuesta para los límites que el proveedor declara y que la respuesta HTTP no refleja.
MailerSend acepta con un \texttt{202} un lote que después rechaza, y Mailchimp Transactional responde con \texttt{200} a una llamada de unos 40 MB cuando declara un tope de 10 MB.
Estos casos plantean qué cuenta como capacidad contractual, y esa cuestión no la resuelve ningún campo nuevo.
La superficie de la API que cambia con el plan ya cabe en v0.3, porque un \textit{endpoint} que un plan no ofrece se omite de ese plan.

Las nueve propuestas y los dieciséis elementos promovidos sin propuesta valen para un dominio y un motor: la recurrencia se cuenta solo en las ocho APIs de correo transaccional, y las Datasheets solo se han evaluado con PRIME4API.

